%% 由 build/emit_tex.py 生成（生成物，勿手改）；定制请放 user_style.tex / user_meta.yaml
\chapter{习题7 参考答案}\label{ux4e60ux98987-ux53c2ux8003ux7b54ux6848}

\begin{quote}
说明：第7章（傅里叶变换）主要习题与参考答案（经典教材整理）。约定
\(\mathcal F[f(t)]=F(\omega)=\int_{-\infty}^{+\infty}f(t)e^{-i\omega t}dt\)。
\end{quote}

\section{习题}\label{ux4e60ux9898-6}

\begin{enumerate}
\def\labelenumi{\arabic{enumi}.}
\item
  求矩形脉冲函数
  \(f(t)=\begin{cases}A,&0\le t\le\tau\\0,&\text{其它}\end{cases}\)
  的傅氏变换。
\item
  求下列函数的傅氏积分：

  \begin{enumerate}
  \def\labelenumii{(\arabic{enumii})}
  \tightlist
  \item
    \(f(t)=\begin{cases}1-|t|,&|t|<1\\0,&|t|\ge1\end{cases}\)；
  \item
    \(f(t)=\begin{cases}1,&|t|<1\\0,&|t|\ge1\end{cases}\)；
  \item
    \(f(t)=e^{-|t|}\)。
  \end{enumerate}
\item
  求下列函数的傅氏变换：

  \begin{enumerate}
  \def\labelenumii{(\arabic{enumii})}
  \tightlist
  \item
    \(f(t)=e^{-a|t|}\ (a>0)\)；
  \item
    \(f(t)=e^{-t^2}\)；
  \item
    \(f(t)=\delta(t-a)\)；
  \item
    \(f(t)=u(t)\)（单位阶跃）。
  \end{enumerate}
\item
  证明傅氏变换的下列性质：

  \begin{enumerate}
  \def\labelenumii{(\arabic{enumii})}
  \tightlist
  \item
    线性性：\(\mathcal F[\alpha f+\beta g]=\alpha F(\omega)+\beta G(\omega)\)；
  \item
    位移性质：\(\mathcal F[f(t-t_0)]=e^{-i\omega t_0}F(\omega)\)；
  \item
    相似性质：\(\mathcal F[f(at)]=\dfrac1{|a|}F\left(\dfrac{\omega}{a}\right)\)；
  \item
    微分性质：\(\mathcal F[f'(t)]=i\omega F(\omega)\)（在适当条件下）。
  \end{enumerate}
\item
  求 \(\delta(t)\) 与 \(\delta(t-t_0)\) 的傅氏变换，并说明其物理意义。
\item
  利用卷积定理证明： \[
  \mathcal F[f*g]=\mathcal F[f]\cdot\mathcal F[g].
  \]
\item
  求下列卷积：

  \begin{enumerate}
  \def\labelenumii{(\arabic{enumii})}
  \tightlist
  \item
    \(u(t)*u(t)\)； (2) \(e^{-t}u(t)*e^{-2t}u(t)\)。
  \end{enumerate}
\item
  利用傅氏变换求微分方程 \(y'(t)+2y(t)=u(t)\) 的解（及稳态解）。
\end{enumerate}

\section{参考答案}\label{ux53c2ux8003ux7b54ux6848-6}

\begin{enumerate}
\def\labelenumi{\arabic{enumi}.}
\item
  \[
  F(\omega)=\int_0^\tau Ae^{-i\omega t}dt=\frac{A}{i\omega}\left(1-e^{-i\omega\tau}\right).
  \]
\item
  \begin{enumerate}
  \def\labelenumii{(\arabic{enumii})}
  \tightlist
  \item
    傅氏积分可按公式计算：\(F(\omega)=\dfrac{2(1-\cos\omega)}{\omega^2}\)（由三角脉冲函数的余弦积分可得）；
  \item
    \(F(\omega)=\dfrac{2\sin\omega}{\omega}\)；
  \item
    \(F(\omega)=\dfrac{2}{1+\omega^2}\)。
  \end{enumerate}
\item
  \begin{enumerate}
  \def\labelenumii{(\arabic{enumii})}
  \tightlist
  \item
    \(F(\omega)=\dfrac{2a}{a^2+\omega^2}\)；
  \item
    \(F(\omega)=\sqrt{\pi}\,e^{-\omega^2/4}\)；
  \item
    \(\mathcal F[\delta(t-a)]=e^{-i\omega a}\)；
  \item
    广义意义下
    \(\mathcal F[u(t)]=\pi\delta(\omega)+\dfrac{1}{i\omega}\)。
  \end{enumerate}
\item
  直接由定义与积分线性性；位移/相似/微分性质均通过换元或分部积分证明。
\item
  \(\mathcal F[\delta(t)]=1\)；\(\mathcal F[\delta(t-t_0)]=e^{-i\omega t_0}\)。物理意义：单位脉冲（瞬时信号）的频谱是``全频带均匀''的。
\item
  \[
  \mathcal F[f*g]=\int_{-\infty}^{+\infty}\left(\int_{-\infty}^{+\infty}f(\tau)g(t-\tau)d\tau\right)e^{-i\omega t}dt
  =\int_{-\infty}^{+\infty}f(\tau)e^{-i\omega\tau}d\tau\int_{-\infty}^{+\infty}g(s)e^{-i\omega s}ds=F(\omega)G(\omega).
  \]
\item
  \begin{enumerate}
  \def\labelenumii{(\arabic{enumii})}
  \tightlist
  \item
    \(u*u=t\,u(t)\)； (2)
    \(e^{-t}u(t)*e^{-2t}u(t)=\left(e^{-t}-e^{-2t}\right)u(t)\)。
  \end{enumerate}
\item
  两边取傅氏变换：\(i\omega Y(\omega)+2Y(\omega)=\pi\delta(\omega)+\dfrac1{i\omega}\)，
  故
  \(Y(\omega)=\dfrac{1}{2+i\omega}\left(\pi\delta(\omega)+\dfrac1{i\omega}\right)\)，经逆变换得
  \(y(t)=\dfrac12\left(1-e^{-2t}\right)u(t)\)。
\end{enumerate}
